% ================= CHƯƠNG 3 =================
\chapter{THIẾT KẾ HỆ THỐNG}\label{chap:3}

Dựa trên cơ sở lý thuyết ở Chương 2, chương này trình bày thiết kế của hệ thống, gồm kiến trúc tổng thể, thiết kế phần cứng các node, thiết kế firmware trên mạng mesh và thiết kế giao tiếp cùng phần mềm giám sát trên máy tính.


\section{Kiến trúc tổng thể}

Hệ thống được tổ chức thành ba lớp như Hình~\ref{fig:3.1}:
\begin{itemize}
    \item \textbf{Lớp mạng mesh:} gồm các node Leaf mang camera và node Root. Leaf chụp ảnh JPEG, chia ảnh thành các chunk và gửi về Root qua mạng mesh; khi cần, một Leaf khác có thể đóng vai trò chuyển tiếp.
    \item \textbf{Lớp liên kết UART:} liên kết tốc độ 921600 bps giữa Root và máy tính, trong đó Root gửi ảnh đã ghép, bản tin trạng thái (telemetry) và log theo một giao thức khung nhị phân có kiểm tra CRC.
    \item \textbf{Lớp phần mềm trên máy tính:} gồm middleware Node.js đọc cổng serial, tách khung và phát dữ liệu qua Socket.IO tới Web Dashboard.
\end{itemize}


% Đặt file figures/h1.png vào project là hình tự hiện
\hinh{h1.png}{0.9}{Sơ đồ khối tổng thể của hệ thống}{fig:3.1}{Chèn hình: sơ đồ khối tổng thể 3 lớp — file baocao/hinh/h1.png trong repo (lưu thành figures/h1.png)}


\section{Thiết kế phần cứng}

Phần cứng của hệ thống gồm một node Root nối với máy tính và các node Leaf mang camera chạy pin độc lập. Mục này trình bày cấu hình của từng loại node, khối nguồn và danh sách linh kiện được sử dụng.


\subsection{Node Root}

Node Root dùng board ESP32-WROOM-32D DevKitC, cấp nguồn và giao tiếp với máy tính qua cổng USB và chip CP2102. Root không xử lý camera nên không cần PSRAM; firmware được điều chỉnh để mọi vùng đệm nằm trong SRAM nội. Thông tin các node trong mạng được tổng hợp trong Bảng~\ref{tab:3.1}.


\begin{table}[!htbp]
\centering
\caption{Vai trò và phần cứng các node}\label{tab:3.1}
{\fontsize{12}{14.5}\selectfont\renewcommand{\arraystretch}{1.25}
\begin{tabular}{|C{1.92cm}|L{2.94cm}|L{9.24cm}|}
\hline
\textbf{Vai trò} & \textbf{Node ID} & \textbf{Phần cứng và ghi chú} \\
\hline
Root & 2846999225 & ESP32-WROOM-32D DevKitC, không PSRAM, CP210x, UART 921600 \\
\hline
Leaf & 1207369673 & ESP32-CAM-MB OV3660, pin 18650 + boost; từng gặp lỗi nguồn \\
\hline
Leaf & 1207433521 & ESP32-CAM-MB OV3660, pin 18650 + boost; chạy ổn định nhất \\
\hline
Leaf & 1207449329 & ESP32-CAM-MB OV3660, pin 18650 + boost \\
\hline
\end{tabular}}
\end{table}


\subsection{Node Leaf và khối nguồn}

Mỗi node Leaf gồm board ESP32-CAM-MB với camera OV3660 và một khối nguồn độc lập để có thể tách node ra xa máy tính. Khối nguồn gồm một pin 18650 loại 2000 mAh, đế pin một cell và module tăng áp DC-DC nhận đầu vào 0,9–5 V, đầu ra 5 V qua cổng USB-A. Đề tài sử dụng mạch tăng áp thay vì mạch hạ áp LM2596 nêu trong tài liệu ban đầu, vì một cell 18650 chỉ có điện áp khoảng 3,7 V, thấp hơn 5 V cần cấp cho board.

Do đầu ra của module là cổng USB-A, board được nối với khối nguồn bằng cáp Micro-USB thông thường; điểm hàn duy nhất là dây từ đế pin vào chân IN của module, kèm một tụ hóa 470 µF hàn song song tại chân IN để giảm sụt áp khi ESP32 phát WiFi với dòng đỉnh lớn. Sơ đồ khối phần cứng một node được trình bày ở Hình~\ref{fig:3.2}.


% Đặt file figures/h2.png vào project là hình tự hiện
\hinh{h2.png}{0.9}{Sơ đồ khối phần cứng node Leaf}{fig:3.2}{Chèn hình: sơ đồ khối phần cứng một node — file baocao/hinh/h2.png trong repo (lưu thành figures/h2.png)}

Trước khi cắm vào board, điện áp đầu ra của module tăng áp được kiểm tra bằng đồng hồ vạn năng theo quy trình hàn tám bước đã soạn. \cansua{Chưa có số liệu đo dòng tiêu thụ của một Leaf khi truyền ảnh và ước tính thời gian chạy pin — không tìm thấy trong tài liệu hiện có, cần đo và bổ sung.}


\subsection{Danh sách linh kiện}

Bảng~\ref{tab:3.2} liệt kê các linh kiện chính đang dùng. Các phụ kiện như tụ hóa, dây cắm và băng keo được mua một lần và dùng chung cho tất cả các node.


\begin{table}[!htbp]
\centering
\caption{Danh sách linh kiện chính}\label{tab:3.2}
{\fontsize{12}{14.5}\selectfont\renewcommand{\arraystretch}{1.25}
\begin{tabular}{|L{6.70cm}|C{1.92cm}|L{5.48cm}|}
\hline
\textbf{Linh kiện} & \textbf{Số lượng} & \textbf{Mục đích} \\
\hline
ESP32-WROOM-32D DevKitC (CP2102) & 1 & Node Root \\
\hline
ESP32-CAM-MB OV3660 (đế nạp CH340) & 3 & Node Leaf/\allowbreak{}Camera \\
\hline
Pin 18650 2000 mAh & 4 & Nguồn Leaf (1 dự phòng) \\
\hline
Đế pin 18650 một cell & 5 & Nguồn Leaf \\
\hline
Module tăng áp DC-DC 0,9–5 V ra 5 V USB & 5 & Nguồn Leaf (dư để dự phòng) \\
\hline
Tụ hóa 470 µF 16 V & 10 & Lọc nguồn tại chân IN của module tăng áp \\
\hline
Dây cắm test board đực – cái 20 cm & 3 bẹ & Đấu nối thử nghiệm \\
\hline
\end{tabular}}
\end{table}


\section{Thiết kế firmware trên mạng mesh}

Firmware trên mạng mesh đảm nhận việc chụp, phân mảnh, truyền và ghép ảnh. Mục này trình bày cấu trúc gói chunk, firmware của node Leaf, firmware của node Root và các tham số cấu hình chính.


\subsection{Cơ chế phân mảnh và cấu trúc gói chunk}

Kích thước một ảnh JPEG lớn hơn nhiều so với giới hạn khoảng 1400 byte của một bản tin painlessMesh, nên ảnh phải được chia thành các mảnh (chunk). Mỗi chunk mang tối đa 1024 byte dữ liệu ảnh và một header 16 byte như Bảng~\ref{tab:3.3}. Header cho phép Root biết chunk thuộc node nào, khung nào, ở vị trí nào và khung có tổng bao nhiêu chunk; trường capture\_us là mốc thời gian chụp theo đồng hồ chung của mesh, được bổ sung ngày 30/09 để đo độ trễ (trước đó header chỉ có 12 byte). Vì painlessMesh truyền bản tin dạng chuỗi, toàn bộ header và payload được mã hóa Base64, nên một gói thực tế dài tối đa 1388 ký tự. Cấu trúc gói chunk trên mạng mesh và khung truyền qua UART được minh họa ở Hình~\ref{fig:3.3}.


\begin{table}[!htbp]
\centering
\caption{Cấu trúc header chunk 16 byte}\label{tab:3.3}
{\fontsize{12}{14.5}\selectfont\renewcommand{\arraystretch}{1.25}
\begin{tabular}{|C{1.92cm}|L{4.65cm}|C{4.14cm}|C{3.10cm}|}
\hline
\textbf{Offset} & \textbf{Trường} & \textbf{Kiểu} & \textbf{Số byte} \\
\hline
0 & node\_\allowbreak{}id & uint32\_\allowbreak{}t & 4 \\
\hline
4 & frame\_\allowbreak{}id & uint16\_\allowbreak{}t & 2 \\
\hline
6 & chunk\_\allowbreak{}index & uint16\_\allowbreak{}t & 2 \\
\hline
8 & total\_\allowbreak{}chunks & uint16\_\allowbreak{}t & 2 \\
\hline
10 & payload\_\allowbreak{}len & uint16\_\allowbreak{}t & 2 \\
\hline
12 & capture\_\allowbreak{}us & uint32\_\allowbreak{}t & 4 \\
\hline
 & Tổng &  & 16 \\
\hline
\end{tabular}}
\end{table}


% Đặt file figures/h3.png vào project là hình tự hiện
\hinh{h3.png}{0.9}{Cấu trúc gói chunk và khung UART}{fig:3.3}{Chèn hình: cấu trúc gói tin chunk và khung UART — file baocao/hinh/h3.png trong repo (lưu thành figures/h3.png)}


\subsection{Firmware node Leaf}

Lưu đồ hàm chụp và gửi ảnh của Leaf được trình bày ở Hình~\ref{fig:3.4}. Theo chu kỳ CAPTURE\_INTERVAL\_MS, Leaf chụp một ảnh, lấy mốc thời gian bằng mesh.getNodeTime(), chia ảnh thành các chunk và gửi lần lượt. Leaf tự bỏ những khung lớn hơn 24 KB để khớp với kích thước vùng đệm của Root.

Ban đầu, Leaf gửi bằng sendBroadcast, khiến mỗi chunk bị phát tới mọi node kể cả các Leaf khác vốn không ghép ảnh. Phiên bản hiện tại để Root phát beacon “\#ROOT” mỗi 2 giây; Leaf học định danh Root từ beacon này và gửi chunk bằng sendSingle thẳng tới Root. Ký tự \# được chọn vì không thuộc bảng mã Base64 nên không bị nhầm với chunk.

Để giảm tranh chấp kênh khi nhiều Leaf cùng gửi, firmware áp dụng cơ chế gửi lệch pha (staggering) dựa trên đồng hồ chung của mesh. Chu kỳ chụp 600 ms được chia thành ba khe 200 ms, mỗi Leaf được gán một khe cố định theo node ID, node không có trong bảng dùng phép chia lấy dư node ID cho số khe.

Leaf còn có cơ chế watchdog: nếu mất beacon của Root quá 8 giây thì quên định danh Root để học lại, quá 30 giây thì khởi động lại board. Mỗi 2 giây, Leaf gửi bản tin \#NSTAT về Root chứa RSSI, heap, phiên bản firmware, nguyên nhân khởi động lại gần nhất và thời gian hoạt động.


% Đặt file figures/h4.png vào project là hình tự hiện
\hinh{h4.png}{0.9}{Lưu đồ chụp và gửi ảnh tại node Leaf}{fig:3.4}{Chèn hình: lưu đồ hàm capture\_and\_send() ở Leaf — file baocao/hinh/h4.png trong repo (lưu thành figures/h4.png)}


\subsection{Firmware node Root}

Root nhận chunk, giải mã Base64 và ghép ảnh theo lưu đồ ở Hình~\ref{fig:3.5}. Mỗi khung đang ghép được giữ trong một vùng đệm (slot) có bitmap đánh dấu các chunk đã nhận và mốc thời gian để hủy khung quá 2 giây chưa đủ chunk. Root có bốn slot, mỗi slot 24 KB, tổng 96 KB SRAM; số slot phải lớn hơn số Leaf gửi đồng thời để chunk của Leaf này không đẩy khung đang ghép dở của Leaf khác ra ngoài.

Khi chunk của một khung mới từ cùng một node tới, Root hủy ngay khung cũ chưa xong của node đó (cơ chế supersede) vì Leaf gửi tuần tự nên khung cũ không thể hoàn thành nữa. Khi hết slot, khung cũ nhất bị loại theo kiểu LRU.

Ngoài chức năng ghép ảnh, Root gửi telemetry mỗi 2 giây gồm số khung ghép thành công và bị hủy, số chunk, heap trống, danh sách node, cây topology lấy từ thư viện, số hop và RSSI từng liên kết, độ trễ trung bình, nhỏ nhất, lớn nhất theo từng node và thời gian tự phục hồi. Root theo dõi thời điểm nhận chunk hợp lệ gần nhất của từng node trong một mảng tĩnh tám phần tử; node có trong danh sách mesh nhưng không gửi chunk quá 10 giây được đánh dấu im lặng để phát hiện hiện tượng node ma. Nếu heap xuống dưới 40 KB liên tục quá 10 giây, Root ghi log và tự khởi động lại.


% Đặt file figures/h5.png vào project là hình tự hiện
\hinh{h5.png}{0.9}{Lưu đồ nhận và ghép ảnh tại node Root}{fig:3.5}{Chèn hình: lưu đồ hàm receive\_and\_reassemble() ở Root — file baocao/hinh/h5.png trong repo (lưu thành figures/h5.png)}


\subsection{Các tham số cấu hình chính}

Các tham số quan trọng của firmware được đặt trong tệp mesh\_config.h của từng node và tổng hợp ở Bảng~\ref{tab:3.4}.


\begin{table}[!htbp]
\centering
\caption{Các tham số cấu hình chính của firmware}\label{tab:3.4}
{\fontsize{12}{14.5}\selectfont\renewcommand{\arraystretch}{1.25}
\begin{tabular}{|L{5.33cm}|L{3.62cm}|L{5.14cm}|}
\hline
\textbf{Tham số} & \textbf{Giá trị} & \textbf{Ý nghĩa} \\
\hline
Payload tối đa mỗi chunk & 1024 byte & Dữ liệu ảnh trong một chunk \\
\hline
MAX\_\allowbreak{}CONCURRENT\_\allowbreak{}FRAMES & 4 & Số khung ghép song song tại Root \\
\hline
MAX\_\allowbreak{}FRAME\_\allowbreak{}SIZE & 24576 byte & Kích thước tối đa một khung \\
\hline
MAX\_\allowbreak{}CHUNKS\_\allowbreak{}PER\_\allowbreak{}FRAME & 24 & Số chunk tối đa một khung \\
\hline
REASSEMBLY\_\allowbreak{}TIMEOUT\_\allowbreak{}MS & 2000 ms & Thời gian chờ ghép một khung \\
\hline
CAPTURE\_\allowbreak{}INTERVAL\_\allowbreak{}MS & 600 ms & Chu kỳ chụp của Leaf (trước là 300 ms) \\
\hline
ROOT\_\allowbreak{}BEACON\_\allowbreak{}INTERVAL\_\allowbreak{}MS & 2000 ms & Chu kỳ phát beacon \#ROOT \\
\hline
ROOT\_\allowbreak{}LOST\_\allowbreak{}RESET\_\allowbreak{}MS /\allowbreak{} RESTART\_\allowbreak{}MS & 8000 /\allowbreak{} 30000 ms & Watchdog Leaf khi mất beacon \\
\hline
STAGGER\_\allowbreak{}SLOTS & 3 & Số khe gửi lệch pha \\
\hline
SILENT\_\allowbreak{}NODE\_\allowbreak{}MS & 10000 ms & Ngưỡng đánh dấu node im lặng \\
\hline
HEAP\_\allowbreak{}CRITICAL\_\allowbreak{}BYTES /\allowbreak{} MS & 40960 byte /\allowbreak{} 10000 ms & Ngưỡng bảo vệ heap tại Root \\
\hline
UART\_\allowbreak{}TX\_\allowbreak{}BUFFER\_\allowbreak{}SIZE & 8192 byte & Bộ đệm phát UART tại Root \\
\hline
\end{tabular}}
\end{table}


\section{Thiết kế giao tiếp và phần mềm trên máy tính}

Mục này trình bày giao thức truyền dữ liệu giữa node Root và máy tính, middleware xử lý dữ liệu trên máy tính và giao diện giám sát Web Dashboard.

\subsection{Giao thức truyền dữ liệu qua UART}

Root gửi dữ liệu lên máy tính qua UART 921600 bps theo giao thức khung nhị phân được đặc tả trong tệp UART\_PROTOCOL.md, cấu trúc như Bảng~\ref{tab:3.5}. Khung bắt đầu bằng hai byte đồng bộ AA 55, tiếp theo là một byte kiểu, bốn byte độ dài, phần dữ liệu và hai byte CRC16/CCITT-FALSE tính trên phần từ kiểu đến hết dữ liệu. Có ba kiểu khung: ảnh đã ghép (0x01), telemetry dạng JSON (0x02) và log (0x03). Khung ảnh có thêm header phụ 10 byte gồm node\_id, frame\_id và độ dài JPEG.


\begin{table}[!htbp]
\centering
\caption{Cấu trúc khung UART}\label{tab:3.5}
{\fontsize{12}{14.5}\selectfont\renewcommand{\arraystretch}{1.25}
\begin{tabular}{|C{3.11cm}|C{2.26cm}|L{8.73cm}|}
\hline
\textbf{Trường} & \textbf{Số byte} & \textbf{Mô tả} \\
\hline
MAGIC & 2 & AA 55, đánh dấu đầu khung \\
\hline
TYPE & 1 & 0x01 ảnh, 0x02 telemetry, 0x03 log \\
\hline
LEN & 4 & Độ dài payload (uint32, little-endian) \\
\hline
PAYLOAD & LEN & Dữ liệu theo kiểu khung \\
\hline
CRC16 & 2 & CRC16/\allowbreak{}CCITT-FALSE trên TYPE..PAYLOAD \\
\hline
\end{tabular}}
\end{table}

Để tránh việc ghi UART chặn vòng lặp chính, Root đặt bộ đệm phát 8192 byte trước khi khởi tạo cổng serial và bỏ khung ảnh nếu dung lượng còn trống của bộ đệm nhỏ hơn một nửa; telemetry và log không bao giờ bị bỏ. Số khung ảnh bị bỏ tại UART được đếm và đưa lên dashboard. Mức log của painlessMesh được hạ xuống chỉ báo lỗi, vì log văn bản chèn vào luồng nhị phân khiến middleware phải bỏ qua nhiều byte để tìm lại đầu khung.


\subsection{Middleware trên máy tính}

Middleware được viết bằng Node.js, đọc cổng serial bằng thư viện serialport và phát dữ liệu tới trình duyệt bằng Socket.IO qua năm sự kiện: connect, disconnect, status, telemetry và frame. Bộ tách khung tìm hai byte MAGIC, kiểm tra độ dài (chấp nhận tối đa 131072 byte) và CRC trước khi phát dữ liệu, như lưu đồ ở Hình~\ref{fig:3.6}. Middleware không ghép ảnh vì Root đã gửi ảnh hoàn chỉnh. Bộ kiểm thử test\_protocol.js đã kiểm tra bảy kịch bản, gồm dữ liệu bị cắt vụn, có byte rác xen giữa và CRC sai, kết quả đạt 7/7.


% Đặt file figures/h6.png vào project là hình tự hiện
\hinh{h6.png}{0.9}{Lưu đồ tách khung tại middleware}{fig:3.6}{Chèn hình: lưu đồ FrameParser.push() ở middleware — file baocao/hinh/h6.png trong repo (lưu thành figures/h6.png)}


\subsection{Web Dashboard}

Dashboard là trang web sử dụng React 18 và Tailwind CSS, nhận dữ liệu thời gian thực từ middleware. Giao diện gồm hai tab chính như Hình~\ref{fig:3.7}:
\begin{itemize}
    \item \textbf{Tab Live:} hiển thị ảnh từ từng drone, trạng thái kết nối, cây topology vẽ bằng SVG với màu cạnh theo RSSI, số hop và độ trễ.
    \item \textbf{Tab So sánh:} tính tốc độ thay đổi của số khung thành công, số khung hủy và số chunk từ telemetry, sau đó nhóm mẫu theo số node để so sánh các chỉ số FPS tổng, FPS mỗi Leaf, thông lượng, tỉ lệ hủy, số chunk mỗi giây, kích thước khung trung bình, heap trống tại Root và thời gian tự phục hồi.
\end{itemize}

Số liệu có thể xuất ra tệp CSV để phân tích. Nguyên tắc thiết kế là không sử dụng dữ liệu giả lập: khi chưa có dữ liệu, dashboard hiển thị trạng thái chưa có dữ liệu.


% Đặt file figures/dashboard.png vào project là hình tự hiện
\hinh{dashboard.png}{0.9}{Giao diện Web Dashboard}{fig:3.7}{Chèn hình: ảnh chụp màn hình Web Dashboard (tab Live và tab So sánh) (lưu thành figures/dashboard.png)}


\section{Kết luận chương}

Chương 3 đã trình bày thiết kế hệ thống theo ba lớp: lớp mạng mesh gồm các node Leaf và node Root, lớp liên kết UART và lớp phần mềm trên máy tính. Ở lớp mạng mesh, ảnh JPEG được phân mảnh thành các chunk tối đa 1024 byte kèm header 16 byte, được gửi thẳng tới Root bằng sendSingle với cơ chế gửi lệch pha, sau đó được ghép tại Root bằng bốn vùng đệm có cơ chế supersede và loại bỏ theo LRU. Ảnh đã ghép được đóng khung nhị phân có CRC16, truyền qua UART tới middleware để tách khung và hiển thị trên Web Dashboard. Chương 4 tiếp theo trình bày kết quả thử nghiệm và đánh giá hệ thống đã thiết kế.
